\chapter*{附件：彩名的临床诊断书}

\noindent%
\htbox{姓名：}~彩名\\
\htbox{年龄：}~24岁\\
\htbox{性别：}~不明\\
\htbox{评估日期：}~2025年7月21日

\halfline%
\subsection*{转介来源与评估原因}

患者系本诊室长期随访个案，本次为常规病情评估及药物管理复诊。由其同事“艾莉丝”陪同（部分时候），并提供部分侧面信息。主要评估目标为症状稳定性、药物依从性及社会功能状况。

\subsection*{主诉}

患者自述：近1月“状态绝佳”，否认任何不适。

陪同者代述：患者近1月日常行为模式相对稳定，能够按时服用处方药物，但在情绪和语言内容方面仍持续存在显著异常。

\subsection*{现病史与既往史}

患者表现出一种长期存在的、系统化的身份认同模式（坚信自己是“世界最后的魔法少女莉莉姆”，肩负着对抗“死灵”的使命）。此模式据信已持续数年，具体起病时间及诱因不详，可能与高强度工作环境压力有关。患者对自身状况缺乏病识感（如，将药物治疗视为维持“魔法力量”的行为）。日常功能在结构化环境及熟悉同伴支持下尚可维持，能够参与工作，但社会适应性明显受限。尿液毒理筛查阴性，排除物质诱发精神病。头颅MRI无异常，排除器质性精神障碍。

\subsection*{精神状况检查（MSE）}

\begin{itemize}
\item\htbox{仪态与行为: }衣着色彩鲜艳、装饰繁复，与其所处日常环境不协调。行为表现活跃，时有夸张的姿态。与评估者互动基本合作，眼神接触良好，语速较快。
\item\htbox{言语与思维：}逻辑连贯。频繁使用特定术语。存在明显且系统化的身份妄想和关系妄想。存在夸大妄想。妄想内容固定，难以通过逻辑辩论动摇。
\item\htbox{感知觉：}目前未报告明确的幻觉体验，但需警惕其将现实刺激进行妄想性解释的可能。
\item\htbox{情绪：}主观报告情绪为“充满活力”、“元气满满”。客观观察情感表达丰富、生动，但时有波动。
\item\htbox{认知功能：}定向力完整。记忆力（尤其是与其妄想体系相关的内容）良好。注意力集中度尚可，但易被其内在“任务”或外部“威胁”分心。抽象思维能力受妄想内容影响，难以进行客观的思考。
\item\htbox{自知力：}完全缺乏。患者不认为自身存在精神疾病，将所有异常体验和行为归因于其身份和使命（“魔法少女”）。
\item\htbox{判断力：}社会判断力受损，表现为不合时宜的着装和言语。但在结构化情境和指导下能完成日常任务（如服药）。
\item\htbox{社会功能：}在无结构或陌生社交场合下功能显著受损，但在熟悉和支持性环境（如，与陪同者相处、参与有明确指导的研究工作）中能维持部分功能。
\end{itemize}

\subsection*{主要诊断}

\begin{itemize}
\item\htbox{F25.0~}\textbf{分裂情感性障碍，双相型，持续性} (Schizoaffective Disorder, Bipolar Type, Continuous)：长期（\(\ge 1\)年）存在的显著精神病性症状（主要是系统化的身份妄想、关系妄想），且独立于心境发作（近1月无心境症状，但精神病性症状仍突出）。同时，存在符合双相障碍标准的情感发作证据（可能是躁狂/轻躁狂或混合状态，表现为持续高能量、欣快情绪、情绪波动剧烈）。精神病性症状和情感症状在疾病过程中占据了主要时间段。
\end{itemize}

\subsection*{鉴别诊断}

\begin{itemize}
\item\htbox{F31.x~}\textbf{双相障碍，伴精神病性特征} (Bipolar Disorder, with Psychotic Features)：患者精神病性症状在无心境发作时仍持续存在，不符合双相障碍伴精神病性特征。
\item\htbox{F20.x~}\textbf{精神分裂症} (Schizophrenia)：若情感症状相对不突出或持续时间远少于精神病性症状，则考虑此诊断，但目前患者情感表现活跃，可能性较低。
\item\htbox{F22~}\textbf{妄想障碍，混合型或夸大型} (Delusional Disorder, Mixed or Grandiose Type)：患者行为怪异性、可能的情感波动及功能损害程度超出典型的妄想障碍。
\end{itemize}

\subsection*{其他需要关注的问题}
\begin{itemize}
\item\htbox{Z60.4~}\textbf{社会排斥或拒绝} (Social Exclusion and Rejection)：可能适用，因患者的怪异行为导致社交困难。
\item\htbox{Z65.8~}\textbf{其他特定社会心理问题} (Other Problems Related to Psychosocial Circumstances)：可能涉及高压力工作环境。
\end{itemize}

\subsection*{评估总结}
患者呈现出典型的、以固定身份妄想为核心的精神病性症状，并伴有显著的情感波动和功能受损。当前最符合的诊断为精神分裂情感性障碍（双相型）。其将现实生活融入妄想体系的行为模式复杂，既维持了部分功能，也固化了病理状态。药物治疗依从性需密切监控。患者缺乏自知力，社会功能受损，需要持续的药物治疗和社会心理支持。

\subsection*{治疗计划}

\begin{itemize}
\item\htbox{药物治疗：}维持现有药物方案，具体药物名称及剂量参见电子病历。密切监测疗效与副作用。强调按时按量服药的重要性，继续利用社会支持系统协助监督。必要时（如出现严重激越或攻击行为）考虑短期辅助镇静药物。
\item\htbox{心理干预：}提供稳定、接纳的治疗关系，不直接挑战其核心妄想，但在安全范围内鼓励现实检验。继续指导陪同者如何与患者有效沟通。视患者接受程度，可尝试引入压力管理、情绪调节及社交技能训练。评估其当前工作环境的压力水平，必要时与相关负责人沟通，寻求合理的职务调整或支持。
\end{itemize}

\subsection*{风险评估}
需关注患者在应激状态下出现冲动、攻击行为或自伤/自杀观念，并评估其妄想内容是否有导致危害他人的风险。

\subsection*{随访}
建议保持每月至少一次的规律复诊，以便及时调整治疗方案。如出现病情显著波动，需随时就诊。

\halfline\hfill\htbox{签字：}~{\large\handwritten Dr.~Toad}
